\section{Reproduction and implementation details}
The GitHub repository separates \texttt{paper/} and \texttt{code/}. The historical research-plan source is preserved under \texttt{paper/research-plan/}; the professor's original template is under \texttt{paper/professor-template-2026/}. Compile this report from \texttt{paper/Final\_Report/} using LuaLaTeX and BibTeX, with \texttt{main.tex} as the entry point.

From \texttt{code/simulation/}, install the project and run the documented test, pilot, evaluation, analysis, benchmark, and replay commands. The frozen design records margins, seed ranges, conditions, and analysis rules. The evaluation manifest records source hashes; raw-log SHA-256 hashes and reconstruction checks accompany the summaries. Bulk raw logs remain local and can be regenerated from the frozen design; a representative raw trace and an offline five-policy replay are supplied for inspection.

The implementation resolves plan details explicitly: observations occur at target/transit cell centres, shortest paths minimize turns with horizontal-first ties, initial heading has no turning cost, slowdown onset is integer-uniform and applies before the tick's movement, and task execution continues within the remaining tick after a waypoint. Failed sessions retain committed ownership; no recovery or theft is introduced. Component operation-count formulas and message encoding are documented beside the code.

The original template contains a stray \texttt{w} in its section-format command. Removing this typographical error preserves the intended formatting. The report retains the template's font family, 11-point body, one-inch margins, spacing, and required headings, and adds GitHub icons and a boxed repository link from the research-plan styling. The KTH logo is omitted in accordance with the course formatting page.
